\begin{frame}[fragile]{Mensagens temporárias com printf}
\begin{lstlisting}
for (size_t i = 0; i < quantidade; ++i) {
    printf("i=%zu valor=%d limite=%d total_antes=%zu\n",
           i, valores[i], limite, total);
    if (valores[i] >= limite) ++total;
    printf("total_depois=%zu\n", total);
}
\end{lstlisting}
\begin{itemize}
\item Registre valores que confirmem ou refutem uma hipótese.
\item \texttt{\%zu} corresponde a \texttt{size\_t}; \texttt{\%d}, a \texttt{int}.
\item Nomeie os campos para interpretar o registro depois.
\item Remova ou desative mensagens temporárias ao concluir.
\end{itemize}
\small Em hardware, mensagens podem mudar a temporização; a observação também tem custo.
\end{frame}
